% =====================================================================
\section{Sistemas de masa variable}
\label{sec:masa-variable}
% =====================================================================

Un cohete que expulsa gas o una gota que capta agua cambian de masa. La
ecuación $\vb{F}=m\vb{a}$ \emph{no} vale para el cuerpo solo; hay que aplicar el
teorema del momento lineal al sistema cuerpo $+$ masa intercambiada.

\begin{teorema}[Ecuación de un cuerpo de masa variable]
\label{teo:masa-variable}
Sea un cuerpo de masa $m(t)$ y velocidad $\vb{v}(t)$ que, en $\dd t$, intercambia
la masa $\abs{\dd m}$ con velocidad $\vb{u}$ \emph{relativa al cuerpo}, sometido a
la fuerza externa $\vb{F}\ext$. Entonces
\begin{equation}
  \resultado{m\,\od{\vb{v}}{t}=\vb{F}\ext+\od{m}{t}\,\vb{u}} .
  \label{eq:masa-variable}
\end{equation}
El término $\dot m\,\vb{u}$ es el \emph{empuje}. Vale tanto si el cuerpo pierde masa
($\dot m<0$, $\vb{u}$ velocidad relativa del gas expulsado) como si la gana
($\dot m>0$, $\vb{u}$ velocidad relativa de la masa captada).
\end{teorema}

\begin{proof}
Caso de expulsión (\cref{fig:cohete-esquema}); el de captación es análogo.
\begin{enumerate}
  \item \emph{Instante $t$}: el sistema es el cuerpo, con $\vb{P}(t)=m\vb{v}$.
  \item \emph{Instante $t+\dd t$}: cuerpo de masa $m+\dd m$ ($\dd m<0$) y velocidad
  $\vb{v}+\dd\vb{v}$, más el gas expulsado, de masa $-\dd m>0$ y velocidad $\vb{v}+\vb{u}$:
  \[
    \vb{P}(t+\dd t)=(m+\dd m)(\vb{v}+\dd\vb{v})+(-\dd m)(\vb{v}+\vb{u}).
  \]
  \item \emph{Balance}. Por el \cref{teo:momento-lineal}, $\dd\vb{P}=\vb{F}\ext\dd t$.
  Desarrollando y despreciando el término de segundo orden $\dd m\,\dd\vb{v}$:
  \[
    \dd\vb{P}=m\,\dd\vb{v}+\vb{v}\,\dd m-\vb{v}\,\dd m-\vb{u}\,\dd m=m\,\dd\vb{v}-\vb{u}\,\dd m=\vb{F}\ext\dd t .
  \]
  \item Dividiendo entre $\dd t$: $m\dot{\vb{v}}=\vb{F}\ext+\dot m\,\vb{u}$.\qedhere
\end{enumerate}
\end{proof}

% ---------------------------------------------------------------------
\begin{figure}[ht]
\centering
\begin{tikzpicture}[scale=1.0]
  % instante t
  \begin{scope}
    \draw[texto!85,line width=0.8pt,fill=caja,rounded corners=2pt] (0,0) rectangle (2.4,0.9);
    \node[etiq] at (1.2,0.45) {$m$};
    \draw[vec,cian] (2.4,0.45) -- ++(1.0,0) node[anchor=west,etiq,inner sep=2pt]{$\vb{v}$};
    \node[etiq,anchor=north,inner sep=6pt] at (1.6,0) {instante $t$};
  \end{scope}
  % instante t+dt
  \begin{scope}[xshift=6.3cm]
    \draw[texto!85,line width=0.8pt,fill=caja,rounded corners=2pt] (0,0) rectangle (2.4,0.9);
    \node[etiq] at (1.2,0.45) {$m+\dd m$};
    \draw[vec,cian] (2.4,0.45) -- ++(1.0,0) node[anchor=west,etiq,inner sep=2pt]{$\vb{v}+\dd\vb{v}$};
    \foreach \p in {(-0.35,0.3),(-0.55,0.62),(-0.75,0.36),(-0.42,0.58)}
      \fill[naranja!80] \p circle (0.06);
    \draw[vec,naranja] (-0.85,0.45) -- ++(-0.9,0) node[anchor=east,etiq,inner sep=2pt]{$\vb{v}+\vb{u}$};
    \node[naranja,font=\footnotesize,anchor=south,inner sep=4pt] at (-0.55,0.75) {$-\dd m$};
    \node[etiq,anchor=north,inner sep=6pt] at (1.2,0) {instante $t+\dd t$};
  \end{scope}
\end{tikzpicture}
\caption{Balance de momento de un cohete. En $\dd t$ expulsa la masa $-\dd m>0$
con velocidad $\vb{u}$ relativa al cohete (en el laboratorio, $\vb{v}+\vb{u}$).}
\label{fig:cohete-esquema}
\end{figure}
% ---------------------------------------------------------------------

\subsection{El cohete}

\begin{notacion}[Transparencias]
\begin{itemize}
  \item $m(t)$: masa del cohete con su combustible; $m_{e}(t)$: masa de gas ya
        expulsada. La masa total se conserva: $m+m_{e}=m_{0}$, luego $\dd m=-\dd m_{e}$.
  \item $v$: velocidad del cohete; $v_{e}$: velocidad (de módulo) del gas en el
        laboratorio, dirigida hacia atrás.
  \item $u\defeq v+v_{e}$: velocidad del gas \emph{relativa al cohete}, que supondremos
        constante ($u\sim2500$--$4500$\,m/s con combustibles químicos, $u\sim10^{5}$\,m/s en
        motores iónicos).
  \item $\alpha\defeq\od{m_{e}}{t}=-\od{m}{t}>0$: ritmo de eyección.
\end{itemize}
\end{notacion}

% ---------------------------------------------------------------------
\begin{figure}[ht]
\centering
\tikzset{
  cohete/.pic={
    % llama
    \shade[top color=ambar,bottom color=naranja!20!fondo,shading angle=0]
      (-0.22,-1.12) .. controls (-0.3,-1.6) and (-0.05,-2.05) .. (0,-2.25)
      .. controls (0.05,-2.05) and (0.3,-1.6) .. (0.22,-1.12) -- cycle;
    \shade[top color=texto,bottom color=ambar!40!fondo]
      (-0.1,-1.12) .. controls (-0.13,-1.4) and (-0.03,-1.65) .. (0,-1.75)
      .. controls (0.03,-1.65) and (0.13,-1.4) .. (0.1,-1.12) -- cycle;
    % aletas
    \fill[violeta!70!caja] (-0.32,-0.55) -- (-0.62,-1.05) -- (-0.32,-0.95) -- cycle;
    \fill[violeta!70!caja] (0.32,-0.55) -- (0.62,-1.05) -- (0.32,-0.95) -- cycle;
    % tobera
    \fill[texto!45] (-0.17,-0.95) -- (0.17,-0.95) -- (0.24,-1.12) -- (-0.24,-1.12) -- cycle;
    % cuerpo
    \shade[left color=texto!55,right color=texto!12,middle color=texto!90]
      (-0.32,-0.95) -- (-0.32,0.75) .. controls (-0.3,1.25) and (-0.1,1.55) .. (0,1.68)
      .. controls (0.1,1.55) and (0.3,1.25) .. (0.32,0.75) -- (0.32,-0.95) -- cycle;
    \draw[texto!80,line width=0.4pt]
      (-0.32,-0.95) -- (-0.32,0.75) .. controls (-0.3,1.25) and (-0.1,1.55) .. (0,1.68)
      .. controls (0.1,1.55) and (0.3,1.25) .. (0.32,0.75) -- (0.32,-0.95) -- cycle;
    % ventanilla
    \fill[cian!70!fondo] (0,0.55) circle (0.11);
    \draw[texto!80,line width=0.4pt] (0,0.55) circle (0.11);
  }
}
% (a) balance
\begin{tikzpicture}[scale=0.95,baseline=(ta.base)]
  \node[etiq,anchor=south] (ta) at (1.8,2.35) {(a) Balance de momento};
  \pic at (0,0) {cohete};
  \draw[vec,cian] (0.75,0.3) -- (0.75,1.3) node[anchor=west,etiq,inner sep=2pt]{$v$};
  \node[etiq,anchor=east,inner sep=6pt] at (-0.32,0.1) {$m$};
  \node[etiq,anchor=north] at (0,-3.3) {instante $t$};
  \begin{scope}[xshift=3.6cm]
    \pic at (0,0.35) {cohete};
    \draw[vec,cian] (0.75,0.65) -- (0.75,1.85) node[anchor=west,etiq,inner sep=2pt]{$v+\dd v$};
    \node[etiq,anchor=east,inner sep=6pt] at (-0.32,0.45) {$m-\dd m_{e}$};
    \foreach \p in {(0.08,-2.3),(-0.12,-2.42),(0.15,-2.55),(-0.05,-2.62)}
      \fill[naranja] \p circle (0.05);
    \draw[vec,naranja] (0.55,-2.25) -- (0.55,-3.0) node[anchor=west,etiq,inner sep=2pt]{$v_{e}$};
    \node[naranja,font=\footnotesize,anchor=east,inner sep=4pt] at (-0.15,-2.45) {$\dd m_{e}$};
    \node[etiq,anchor=north] at (0,-3.3) {instante $t+\dd t$};
  \end{scope}
\end{tikzpicture}\hspace{14mm}%
% (b) fuerzas
\begin{tikzpicture}[scale=0.95,baseline=(tb.base)]
  \node[etiq,anchor=south] (tb) at (0,2.35) {(b) Fuerzas en el ascenso};
  \pic at (0,0) {cohete};
  \draw[eje] (-2.0,-3.0) -- (-2.0,1.9) node[anchor=south,peq,inner sep=2pt]{$z$};
  \draw[vec,verde] (0.75,-0.4) -- (0.75,1.1) node[anchor=west,etiq,inner sep=2pt]{$F_{\text{emp}}=\alpha u$};
  \draw[vec,rosa] (-0.85,-0.1) -- (-0.85,-1.2) node[anchor=east,etiq,inner sep=2pt]{$mg$};
  \draw[vec,ambar] (-0.85,1.35) -- (-0.85,0.45) node[anchor=east,etiq,inner sep=3pt,pos=0.5]{$F_{\text{roz}}$};
  \draw[vec,naranja] (0.5,-2.1) -- (0.5,-2.85) node[anchor=west,etiq,inner sep=2pt]{$-u\,\kk$};
\end{tikzpicture}
\caption{Cohete. (a) Balance de momento entre $t$ y $t+\dd t$: el cohete
pierde $\dd m_{e}$, que sale hacia atrás con velocidad $v_{e}$ en el laboratorio.
(b) Fuerzas en el ascenso vertical: peso, rozamiento opuesto a $v$ y empuje
$\alpha u$, reacción al gas expulsado con velocidad $-u\,\kk$ relativa al cohete.}
\label{fig:cohete}
\end{figure}
% ---------------------------------------------------------------------

\begin{proposicion}[Ecuación del cohete libre]
\label{prop:cohete-libre}
Sin fuerzas externas, la conservación del momento lineal entre $t$ y
$t+\dd t$ (\cref{fig:cohete}a), a primer orden, da
\begin{gather}
  m v=(m-\dd m_{e})(v+\dd v)-\dd m_{e}\,v_{e}
  \quad\Longrightarrow\quad
  m\,\dd v-\dd m_{e}\,(v+v_{e})=0,\notag\\
  \resultado{m\,\dd v=u\,\dd m_{e}=-u\,\dd m} .
\end{gather}
\end{proposicion}

\begin{proof}
Desarrollando, $(m-\dd m_{e})(v+\dd v)=mv+m\,\dd v-v\,\dd m_{e}-\dd m_{e}\,\dd v$; se
desprecia $\dd m_{e}\,\dd v$ (segundo orden) y se agrupa con $-v_{e}\,\dd m_{e}$. Por último,
$v+v_{e}=u$ y $\dd m_{e}=-\dd m$. Es el \cref{teo:masa-variable} con $\vb{u}=-u\,\kk$ y $\dot m=-\alpha$.
\end{proof}

\begin{corolario}[Tsiolkovski]
\label{cor:tsiolkovski}
Partiendo de $v_{0}$ con masa $m_{0}$,
\begin{equation}
  \int_{v_{0}}^{v}\dd v'=-u\int_{m_{0}}^{m}\frac{\dd m'}{m'}
  \quad\Longrightarrow\quad
  \resultado{v-v_{0}=u\ln\frac{m_{0}}{m}}
  \quad\Longleftrightarrow\quad
  m=m_{0}\,\e^{-(v-v_{0})/u} .
  \label{eq:tsiolkovski}
\end{equation}
Si $v_{0}=0$, la masa expulsada es $m_{e}=m_{0}-m=m_{0}\bigl(1-\e^{-v/u}\bigr)$. La velocidad depende
sólo del cociente de masas, no del ritmo de consumo.
\end{corolario}

\begin{corolario}[Eyección a ritmo constante]
\label{cor:cohete-alfa}
Si $\alpha$ es constante, $m=m_{0}-\alpha(t-t_{0})$ y, con $v_{0}=0$,
\begin{equation}
  \resultado{v(t)=-u\ln\Bigl[1-\frac{\alpha}{m_{0}}(t-t_{0})\Bigr]} .
\end{equation}
\end{corolario}

\begin{proposicion}[Ascenso vertical en un campo uniforme]
\label{prop:cohete-g}
Con $\vb{F}\ext=-mg\,\kk$, la ecuación del cohete es
\begin{equation}
  \vb{F}\ext=-mg\,\kk=\Bigl(u\od{m}{t}+m\od{v}{t}\Bigr)\kk
  \qquad\Longrightarrow\qquad
  m\od{v}{t}=\alpha u-mg ,
\end{equation}
y con $\alpha$ constante, $t_{0}=0$ y $v_{0}=0$,
\begin{equation}
  \resultado{v(t)=-u\ln\Bigl(1-\frac{\alpha t}{m_{0}}\Bigr)-gt}=u\ln\frac{m_{0}}{m(t)}-gt .
\end{equation}
\begin{itemize}
  \item \emph{Empuje}: es opuesto al flujo de momento de los gases,
        $\vb{F}_{\text{emp}}=-\od{m_{e}}{t}\,\vb{u}_{\text{gas}}=\alpha u\,\kk$, con $\vb{u}_{\text{gas}}=-u\,\kk$.
  \item \emph{Despegue}: el cohete acelera hacia arriba ($\dot v>0$) si
        $F_{\text{emp}}>mg$, es decir, $\resultado{\alpha u>mg}$.
  \item \emph{Rozamiento}: en la atmósfera se añade
        $\vb{F}_{\text{roz}}=-\tfrac12\rho C_{D}A\,v\abs{v}\,\kk$, opuesto a la velocidad; al ascender,
        $v>0$ y $v\abs{v}=v^{2}$.
\end{itemize}
\end{proposicion}

\begin{proof}
$\dot v=\alpha u/(m_{0}-\alpha t)-g$; integrando de $0$ a $t$,
$v=-u\bigl[\ln(m_{0}-\alpha t)-\ln m_{0}\bigr]-gt$.
\end{proof}

% ---------------------------------------------------------------------
\begin{figure}[ht]
\centering
\begin{tikzpicture}
\begin{groupplot}[group style={group size=2 by 1,horizontal sep=14mm},
  oscuro,width=7.4cm,height=6cm]
  \nextgroupplot[xmin=0,xmax=0.92,ymin=0,ymax=2.6,
    xlabel={$\alpha t/m_{0}$},ylabel={$v/u$},
    title={(a) Velocidad},legend pos=north west]
    \addplot[cian,domain=0:0.9,samples=150,name path=c0] {-ln(1-x)};
    \addlegendentry{$\gamma=0$}
    \addplot[ambar,domain=0:0.9,samples=150] {-ln(1-x)-0.5*x};
    \addlegendentry{$\gamma=0{,}5$}
    \addplot[rosa,domain=0:0.9,samples=150,name path=c9] {-ln(1-x)-0.9*x};
    \addlegendentry{$\gamma=0{,}9$}
    \addplot[cian,opacity=0.10] fill between[of=c0 and c9];
  \nextgroupplot[xmin=0,xmax=3.05,ymin=0,ymax=1.05,
    xlabel={$\Delta v/u$},ylabel={$m/m_{0}$},
    title={(b) Fracción de masa restante}]
    \addplot[violeta,domain=0:3,samples=150,name path=m] {exp(-x)};
    \path[name path=ax] (axis cs:0,0) -- (axis cs:3,0);
    \addplot[violeta,opacity=0.12] fill between[of=m and ax];
    \draw[aux] (axis cs:1.756,0) -- (axis cs:1.756,0.1727) -- (axis cs:0,0.1727);
    \node[punto,fill=ambar,inner sep=1.5pt] at (axis cs:1.756,0.1727) {};
    \node[ambar,font=\footnotesize,anchor=south west,inner sep=3pt,align=left] at (axis cs:1.756,0.1727)
      {órbita baja:\\ $\Delta v\simeq7{,}9$\,km/s,\\ $u=4{,}5$\,km/s};
\end{groupplot}
\end{tikzpicture}
\caption{(a) Velocidad de un cohete con eyección constante, para
$\gamma\defeq m_{0}g/(\alpha u)$; la franja muestra lo que «roba» la gravedad
($\gamma\,\alpha t/m_{0}$). (b) Tsiolkovski, $m/m_{0}=\e^{-\Delta v/u}$: para ponerse en
órbita baja con propulsión química sólo llega arriba el 17\,\% de la masa
inicial; por eso los cohetes son casi todo combustible y van por etapas.}
\label{fig:cohete-v}
\end{figure}
% ---------------------------------------------------------------------

\subsection{Captación de masa en reposo}

Si el cuerpo capta masa que estaba \emph{en reposo} en el laboratorio
(una gota en una nube), la velocidad relativa de la masa captada es
$\vb{u}=\vb{0}-\vb{v}=-\vb{v}$, y \eqref{eq:masa-variable} queda
\begin{equation}
  m\dot{\vb{v}}=\vb{F}\ext-\dot m\,\vb{v}
  \qquad\Longleftrightarrow\qquad
  \resultado{\od{(m\vb{v})}{t}=\vb{F}\ext} .
  \label{eq:captacion}
\end{equation}
En este caso, y sólo en él, vale «$\vb{F}=\dot{\vb{p}}$» para el cuerpo (\cref{prob:8}).
